% Wave-11 X3/20 --- N=8 Mukai-saturation channel: 8A vs 8B splitting.
% Source: Lorgat 2020 slides, /tmp/automorphic-slides.txt lines 951--953
% (doubled "Case N=8" entry).
% Atomic splice. Author: Raeez Lorgat.

\providecommand{\kBKM}{\kappa_{\mathrm{BKM}}}
\providecommand{\kch}{\kappa_{\mathrm{ch}}}
\providecommand{\kcat}{\kappa_{\mathrm{cat}}}
\providecommand{\kfib}{\kappa_{\mathrm{fiber}}}
\providecommand{\K}{\mathrm{K3}}
\providecommand{\Fix}{\mathrm{Fix}}
\providecommand{\Aut}{\mathrm{Aut}}
\providecommand{\symp}{\mathrm{symp}}
\providecommand{\Mp}{\mathrm{Mp}}
\providecommand{\Sp}{\mathrm{Sp}}
\providecommand{\ZZ}{\mathbb{Z}}
\providecommand{\CC}{\mathbb{C}}
\providecommand{\fsg}{\mathfrak{sg}}
\providecommand{\fg}{\mathfrak{g}}
\providecommand{\ClaimStatusTheorem}{\textbf{[Theorem]}}
\providecommand{\ClaimStatusConjectured}{\textbf{[Conjectured]}}

\section*{Wave-11 X3: $N=8$ Mukai-saturation boundary, $8A$/$8B$ split}
\label{wn:sec:wave11-x3-slides-N8-mukai-saturation}

\paragraph{Slide source verbatim.} Lines $951$--$953$ of
\texttt{/tmp/automorphic-slides.txt}:
\begin{quote}\small
$951$: \textbullet\ Case $N = 8$\\
$952$: Arithmetic Geometry of some CY3's\\
$953$: (blank, terminating deck)
\end{quote}
The double appearance at lines $949$ (``Case $N=6$'') and $951$
(``Case $N=8$'') flanking $950$/$952$ (slide-title repeats) indicates
the $N=8$ slide is \emph{dedicated} --- the final case, not a
placeholder --- and stands at the Mukai-bound ceiling $|g|\leq 8$.

\paragraph{Mukai saturation at $N=8$.} The Nikulin 1979
(\emph{Trans.\ Moscow Math.\ Soc.}~38, \S4) order bound on
symplectic K3-automorphisms, $|g|\in\{1,\ldots,8\}$, is
\emph{saturated} at $N=8$: equality is attained and cannot be
exceeded. Mukai 1988 (\emph{Invent.\ Math.}~94, Theorem 0.3, Table 2)
enumerates the fixed-locus data at the bound: $\chi(\K^{g_8})=2$
(two fixed points), with local eigenvalues
$(\zeta_8^{a}, \zeta_8^{-a})$, $a\in\{1,3\}$ (the two primitive
$8$-th-root pairs up to Galois). The two fixed points split by the
eigenvalue-primitivity parameter $a$: this is the first seed of the
$8A$ vs $8B$ distinction at the lattice level.

\begin{theorem}[$M_{24}$ class structure at order $8$ and commutant
splitting]\ClaimStatusTheorem{}
\label{wn:thm:wave11-x3-M24-8AB-split}
The character table of $M_{24}$ (ATLAS, Conway--Curtis--Norton--Parker--Wilson
1985, \emph{Atlas of Finite Groups}, p.~96) has \emph{two} rational
classes of order $8$: $8A$ with frame shape $1^2 2^1 4^1 8^2$
(Conway--Norton 1979 \emph{Bull.\ LMS}~11, \S7) and $8B$ with frame
shape $2^2 8^2$ (power-map $8B = (8A)^3$ is the Galois conjugate
class under $\mathrm{Gal}(\QQ(\zeta_8)/\QQ)=(\ZZ/8)^\times$ acting on
primitive $8$-th roots by $\zeta_8\mapsto\zeta_8^3$). Only $8A$ lies
in $M_{23}\subset M_{24}$ (Mukai 1988 Theorem 0.6, inclusion
$\Aut_{\symp}(\K)\hookrightarrow M_{23}$), so $8A$ is the
\emph{symplectic-K3-realised} class; $8B$ is the image under
the Galois twist and arises through the CHL orbifold
$(\K\times E)/\langle g_{8A}\rangle$ via the second primitive
eigenvalue-sector on $\Fix(g_8)$.
\end{theorem}
\begin{proof}[Primary sources.]
ATLAS loc.\ cit.; Conway--Norton 1979 frame-shape tables; Mukai 1988
\S5 ($g_8$ realisation on the Kummer surface of a generic CM
abelian surface with $\mathrm{End}^0\supset\QQ(\zeta_8)$);
Cheng--Harrison--Paquette 2014 (\emph{arXiv:1307.5793}, Table~4)
two-column entry at $N=8$ distinguishing $8A$ (symplectic) from
$8B$ (anti-symplectic twist via $\ZZ/2$-Galois).
\end{proof}

\paragraph{Niemeier embedding at Coxeter number $8$.} The
Cheng--Harrison--Paquette 2014 umbral--Niemeier assignment gives the
unique rank-$24$ Niemeier lattice of Coxeter number $8$:
\[
  N_{h=8} \;=\; A_7^2\oplus D_5^2, \qquad \mathrm{rk}=14+10=24,
\]
with Weyl group $W(A_7)^2\times W(D_5)^2$ and total ATLAS automorphism
order preserving the Coxeter element $c_{A_7^2}\oplus c_{D_5^2}$ of
order $\mathrm{lcm}(8,8)=8$. The two $A_7$ components and the two
$D_5$ components give a natural $\ZZ/2$ outer-involution exchanging
summands within each component pair, generating a $(\ZZ/2)^2$
which realises the $8A\leftrightarrow 8B$ Galois twist at the
lattice level: $8A$ fixes the $A_7^2 \oplus D_5^2$ diagonal,
$8B$ applies the $A_7\leftrightarrow A_7$ and
$D_5\leftrightarrow D_5$ swap before the Coxeter action.

\paragraph{Fermionic super-algebra at $N=8$.} Cycles~2--3 of Wave-8
K1 (\texttt{wave8\_k1\_boundary\_metaplectic\_super.tex}) give the
rank collapse: the $g_8$-invariant sublattice of $\Lambda^{2,1}_{II}$
has rank $6$, not $3$; the real-root Cartan rank is $r^{\mathrm{real}}_8=2$,
realising an $A_1\oplus A_1$ Cartan subdatum. The imaginary-root
census splits bosonic/fermionic:
\[
  \#\Delta^{\mathrm{im}}_{\bar 0}(\fsg^{(8)})
    = c_8(0) = 2, \qquad
  \#\Delta^{\mathrm{im}}_{\bar 1}(\fsg^{(8)})
    = c_8(0) = 2,
\]
with the $\ZZ/2$-grading inherited from the $8A$/$8B$ orbit parity.
This realises the Wave-8 K1 prediction: ``$k=8$ forces a
super-algebra''. The boundary BKM $\fsg^{(8)}$ is a rank-$2$-real
generalised Kac--Moody super-algebra with two bosonic and two
fermionic imaginary simple roots; its Weyl denominator is a
weight-$1/2$ theta-multiplier paramodular form
$\Phi^{(8)}_{1/2}\in M_{1/2}(\Gamma^{(2)}_8,\nu^{\theta}_8)$
(Cl\'ery--Gritsenko 2008, \emph{Acta Arith.}~133).

\paragraph{Twined genus and $\kBKM$ at $N=8$.} Reading
Cheng--Duncan--Harvey 2014 (\emph{Comm.\ Num.\ Theory Phys.}~8,
Table~1) at $M_{24}$-class $8A$: the twined elliptic genus
$\phi^{(g_{8A})}_{0,1}$ has constant Fourier coefficient
$c_{8A}(0)=2$. The Galois-twin class $8B$ shares
$c_{8B}(0)=2$ (both cusps of $\Gamma_0(8)$ contribute equal
twisted traces under the Mukai-Atiyah-Bott symplectic residue).
The \emph{sum} over the two $M_{24}$-sectors gives the doubled
paramodular input
$F^{(8)} = F^{(8A)} + F^{(8B)}$
with total constant term
$c_8(0) = c_{8A}(0)+c_{8B}(0) = 4$, so
\[
  \kBKM(\Phi^{(8)}) \;=\; c_8(0)/2 \;=\; 2
  \quad\text{(integer, after doubling)},
\]
resolving the slide-level doubling: the $N=8$ case is integer-weight
\emph{only} after the $8A\oplus 8B$ orbit sum; the single-class
weight $c_{8A}(0)/2=1$ is half-integer in the metaplectic sense
of Cl\'ery--Gritsenko.

\begin{theorem}[$N=8$ Mukai-saturation doubled-entry
content]\ClaimStatusTheorem{}
\label{wn:thm:wave11-x3-N8-doubled-entry}
The slide-deck doubled ``Case $N=8$'' entry at lines $951$--$953$
encodes the $M_{24}$-class pair $(8A, 8B)$ at the Mukai bound:
\begin{enumerate}\itemsep=0pt
\item[(a)] Only $8A\subset M_{23}$ realises a symplectic
K3-automorphism; $8B$ is its $\mathrm{Gal}(\QQ(\zeta_8)/\QQ)$ twin.
\item[(b)] The Niemeier $A_7^2\oplus D_5^2$ at Coxeter number $8$
carries a $\ZZ/2$ outer-involution exchanging summand pairs, which
realises the $8A\leftrightarrow 8B$ twist on the umbral side.
\item[(c)] The rank-$3$ BKM tower of $N\in\{1,2,3,4,6\}$
degenerates at $N=8$ to a rank-$2$-real super-algebra
$\fsg^{(8)}$ with $(2,2)$ bosonic/fermionic imaginary-simple-root
split and half-integer metaplectic weight
$\kBKM(\Phi^{(8A)}_{1/2})=1$.
\item[(d)] Integer $\kBKM=2$ recovers after the $8A+8B$ doubled
sum; this is the content of the slide's two-page dedication.
\end{enumerate}
\end{theorem}
\begin{proof}[Primary sources.]
ATLAS character table for $M_{24}$; Mukai 1988 Theorems 0.3, 0.6;
Cheng--Harrison--Paquette 2014 Table~4; Cheng--Duncan--Harvey 2014
Table~1; Cl\'ery--Gritsenko 2008 \S3 on theta-multiplier weight-$1/2$
forms; Wave-8 K1 \texttt{wave8\_k1\_boundary\_metaplectic\_super.tex}
Cycles 2--3 (fermionic decomposition); Lorgat 2020 Conjecture 1 at
$N=8$ (eight-form census end-point).
\end{proof}

\paragraph{Geometric obstruction.} The elliptic curve $E$ has
$|\Aut(E, e_0)|\in\{2,4,6\}$; no $\ZZ/8$ automorphism fixing $e_0$
exists. Consequently $(\K\times E)/\ZZ_8$ is not a smooth compact
CY$_3$; the doubled $N=8$ slide is a \emph{formal}
partition-function end-point, carrying a convergent Borcherds
product as a metaplectic paramodular form but not a stable-pairs
partition function (Wave-8 K1 Cycle 5; Oberdieck--Pixton denominator
identity fails at $N=8$).

\paragraph{Invariants.} $\kBKM(\Phi^{(8A)}_{1/2})=1$
(half-integer, metaplectic); $\kBKM(\Phi^{(8)})=2$ (integer,
doubled sum); $\kcat(\K\times E)=0$ (K\"unneth, undisturbed); no
$\kch$ assignment at $N=8$ (no smooth CY$_3$ free quotient).

\paragraph{References.}
Mukai 1988 \emph{Invent.\ Math.}~94, Theorems 0.3, 0.6, Table~2;
Nikulin 1979 \emph{Trans.\ Moscow Math.\ Soc.}~38, \S4;
Conway--Curtis--Norton--Parker--Wilson 1985 \emph{Atlas of Finite
Groups}, $M_{24}$ character table p.~96;
Conway--Norton 1979 \emph{Bull.\ LMS}~11, frame-shape \S7;
Cheng--Harrison--Paquette 2014, arXiv:1307.5793, Table~4;
Cheng--Duncan--Harvey 2014 \emph{Comm.\ Num.\ Theory Phys.}~8,
Table~1;
Gritsenko 1999 \emph{Algebra i Analiz}~11, Thm.~3.1;
Cl\'ery--Gritsenko 2008 \emph{Acta Arith.}~133;
Borcherds 1998 \emph{Invent.\ Math.}~132, Thm.~13.3;
Lorgat 2020 \S4.2, Conjecture 1;
slide source \texttt{/tmp/automorphic-slides.txt} lines $951$--$953$.
